% =====================================================================
% TOHSENO — Genesis Whitepaper
% Typeset for permanence. Compile with XeLaTeX.
%   latexmk -xelatex tohseno-whitepaper.tex
% =====================================================================
\documentclass[11pt]{article}

% ---------- geometry & base ----------
\usepackage[letterpaper,
            textwidth=5.4in,
            textheight=8.6in,
            hcentering,
            top=1.1in]{geometry}
\usepackage{fontspec}
\usepackage{microtype}
\usepackage{xcolor}
\usepackage{graphicx}
\usepackage{fancyvrb}
\usepackage{fancyhdr}
\usepackage{titlesec}
\usepackage{titletoc}
\usepackage{enumitem}
\usepackage{needspace}
\usepackage[hidelinks]{hyperref}

% ---------- the darkroom palette ----------
\definecolor{ink}{HTML}{11100E}       % near-black
\definecolor{halide}{HTML}{77736B}    % silver-halide grey
\definecolor{faint}{HTML}{AAA69D}     % pale grey
\definecolor{signal}{HTML}{FF9D18}    % the safelight
\definecolor{signaldeep}{HTML}{8C4D00}% safelight, printable

\color{ink}

% ---------- fonts ----------
\setmainfont{EBGaramond-Regular.otf}[
  ItalicFont     = EBGaramond-Italic.otf,
  BoldFont       = EBGaramond-Bold.otf,
  BoldItalicFont = EBGaramond-BoldItalic.otf,
  Numbers = OldStyle,
  Ligatures = TeX
]
\setmonofont{IBMPlexMono-Regular.otf}[
  ItalicFont     = IBMPlexMono-Italic.otf,
  BoldFont       = IBMPlexMono-SemiBold.otf,
  Scale=0.78
]
\newfontfamily\displayfont{EBGaramond-Medium.otf}[
  ItalicFont = EBGaramond-MediumItalic.otf,
  Ligatures  = TeX
]
\newfontfamily\monolabel{IBMPlexMono-Medium.otf}[Scale=0.72, LetterSpace=16.0]

% ---------- page numbers as exposures: 001, 002, ... ----------
\newcommand{\shotpage}{%
  \ifnum\value{page}<10 00\arabic{page}%
  \else\ifnum\value{page}<100 0\arabic{page}%
  \else\arabic{page}\fi\fi}

\pagestyle{fancy}
\fancyhf{}
\renewcommand{\headrulewidth}{0pt}
\fancyhead[C]{{\monolabel\scriptsize\color{halide}TOHSENO / GENESIS WHITEPAPER}}
\fancyfoot[C]{{\ttfamily\scriptsize\color{halide}\shotpage}}
\fancypagestyle{plain}{%
  \fancyhf{}%
  \renewcommand{\headrulewidth}{0pt}%
  \fancyfoot[C]{{\ttfamily\scriptsize\color{halide}\shotpage}}%
}

% ---------- section styling ----------
\titleformat{\section}
  {\needspace{6\baselineskip}\displayfont\LARGE}
  {{\ttfamily\large\color{signaldeep}\thesection}}{0.9em}{}
\titlespacing*{\section}{0pt}{2.6em}{1.1em}

\titleformat{\subsection}
  {\displayfont\large\bfseries}
  {{\ttfamily\normalsize\color{signaldeep}\thesubsection}}{0.8em}{}
\titlespacing*{\subsection}{0pt}{1.8em}{0.6em}

% ---------- verbatim: protocol blocks ----------
\DefineVerbatimEnvironment{proto}{Verbatim}
  {fontsize=\small, frame=leftline, rulecolor=\color{signal},
   framerule=1.2pt, framesep=4mm, xleftmargin=1mm,
   baselinestretch=1.05}

% ---------- pull-quote / creed ----------
\newenvironment{creed}
  {\begin{center}\begin{minipage}{0.86\textwidth}\itshape\large\centering}
  {\end{minipage}\end{center}\vspace{0.4em}}

\newcommand{\creedrule}{\begin{center}{\color{signal}\rule{2.2em}{1.2pt}}\end{center}}

% ---------- honesty-vocabulary tag ----------
\newcommand{\statustag}[1]{{\monolabel\scriptsize\color{signaldeep}#1}}

% ---------- misc ----------
\setlength{\parskip}{0.45em}
\setlength{\parindent}{0pt}
\emergencystretch=1.6em
\newcommand{\arrowdown}{\begin{center}$\downarrow$\end{center}\vspace{-0.6em}}
\setlist{itemsep=0.15em, topsep=0.3em}
\newcommand{\TOH}{TOHSENO}

\begin{document}

% =====================================================================
% TITLE PAGE
% =====================================================================
\begin{titlepage}
\thispagestyle{empty}
\vspace*{1.6in}
\begin{center}

{\displayfont\fontsize{46}{50}\selectfont TOHSENO\par}
\vspace{0.05in}
{\color{faint}\scalebox{1}[-0.72]{\displayfont\fontsize{46}{50}\selectfont TOHSENO}\par}

\vspace{0.65in}
{\displayfont\Large A Protocol for Turning Intention into Living Software\par}

\vspace{0.45in}
{\monolabel\small GENESIS WHITEPAPER AND CONSTITUTIONAL SPECIFICATION\par}
\vspace{0.7em}
{\ttfamily\small Version 0.2 --- July 26, 2026\par}
{\ttfamily\small Status: Foundational, pre-stable-release (factory 0.5.0)\par}
{\ttfamily\small Origin: JP Franetovic / Anky, Inc.\par}

\vspace{0.9in}
\begin{minipage}{0.7\textwidth}
\centering\itshape
Give every idea a shot.\\[0.5em]
The factory is local. Identity is held. Memory is networked.\\[0.5em]
Most shots miss. The prototype is the payoff.
\end{minipage}

\vfill
{\monolabel\scriptsize\color{halide}TOH / ROLL 01 / WRITTEN BEFORE THE FIRST STABLE RELEASE\par}

\end{center}
\end{titlepage}

% =====================================================================
% A NOTE TO WHOEVER FINDS THIS
% =====================================================================
\section*{A note to whoever finds this}

\TOH{} began with a simple observation:

\begin{creed}
A person can carry a coherent piece of software in their mind long before
they know how to build it, distribute it, maintain it, or explain it to
anyone else.
\end{creed}

For most of the history of computing, converting that intention into
software required a project: a team, a roadmap, a repository, a product
manager, infrastructure, accounts, capital, and months of sustained
coordination. The cost of producing software was so high that only a small
fraction of human intentions were allowed to become executable.

That assumption is breaking.

The cost of expressing software is collapsing. A sufficiently clear
intention can increasingly become a functioning application in a single
sustained act. But the systems around software have not caught up. Creation
may become fast while everything that gives an application a life remains
fragmented: ownership, evolution, testing, feedback, identity,
distribution, reputation, provenance, monetization, and memory.

\TOH{} exists to close that circle.

It is not merely a generator. It is not a cloud IDE, a token launcher, a
marketplace, or a wrapper around a coding agent. It is an open protocol and
local-first factory for turning a coherent intention into an independent
application, allowing that application to evolve through contact with
reality, and optionally connecting it to a shared network without
surrendering its sovereignty.

And it makes a second observation, quieter but just as load-bearing:

\begin{creed}
Most shots miss.\\
That is not a tragedy to be engineered away.\\
It is the reason the factory must make taking another one cheap.
\end{creed}

A photographer does not mourn a frame that did not land. The frame goes on
the contact sheet, it teaches the eye, and the roll continues. \TOH{}
treats software the same way. An app idea is a Shot, not a forever bet.
Every Shot lands in the builder's contact sheet as a working artifact they
can use, change, keep, or move past. The prototype is the payoff. Volume
reveals signal. Taste emerges through making.

This document records the system before its first stable release.

It is written so that a future builder, machine, node operator, historian,
or stranger can reconstruct not only how \TOH{} worked, but what it
believed software should become.

\creedrule

% =====================================================================
% ABSTRACT
% =====================================================================
\section*{Abstract}

\TOH{} introduces the \textbf{Shot} as a primitive for personal software.

A Shot is one coherent intention embodied as an independent application and
carried forward through a continuous lineage. The first intent creates the
Shot. Subsequent intents do not create new Shots; they create
\textbf{Evolutions} of the same Shot.

A Shot has three canonical lifecycle states:

\begin{proto}
EVOLVING -> PUBLISHED -> APP_STORE
\end{proto}

\textbf{EVOLVING} means the Shot exists primarily within the builder's
local environment. \textbf{PUBLISHED} means its source is available through
the \TOH{} ecosystem so that other people can inspect, download, build,
run, and test it. \textbf{APP\_STORE} means the Shot has been released
through Apple's App Store. These states are distribution milestones, not
the end of development. A Shot may continue evolving indefinitely after
publication or release.

The \TOH{} factory is open-source and local-first. A builder must be able
to install \TOH{}, create a Shot, evolve it, verify it, publish its source,
and ship it without a \TOH{} account, mobile app, wallet, blockchain
transaction, token, or official server. Every generated application must
remain independent and operable if \TOH{} disappears --- \emph{ejectable
from birth}, a property true at creation time, not a later export feature.

Inside the factory, responsibility is divided by a strict boundary:
artificial intelligence interprets meaning; deterministic code owns
composition, locks, storage invariants, verification, and authority. The
coding agent builds, but \TOH{} --- deterministic, verifiable, unimpressed
--- is the only party allowed to say the result works.

Around that sovereign local factory, \TOH{} defines an optional continuity
layer: portable builder identity; app-specific user identity; explicit
cross-app consent; mobile signing and approvals; public builder records;
structured feedback; sharpened next intents; discovery; node-backed
availability; append-only provenance; and deployment-agnostic appcoin
links.

The future \TOH{} mobile app is a premium client for this continuity
layer. It is not mandatory and it is not the root dependency of generated
applications. It is subject to a genesis invariant: \textbf{it must not be
built by hand before the protocol is stable. It must be the first canonical
Shot created by the first stable \TOH{} release.}

\TOH{} nodes are compatible services that preserve and serve public
memory. Nodes index; builders own. The official node is only the first
node, not the definition of the network. An append-only registry may anchor
public claims on an existing blockchain; it records small signed facts and
content hashes, and nothing else. Appcoin deployment remains agnostic:
\TOH{} records the relationship between a Shot and its canonical appcoin,
and does not deploy, mint, trade, price, or endorse the asset.

The result is a system in which:

\begin{proto}
one intention becomes one Shot;
one Shot can evolve for years;
the source remains owned by the builder;
identity remains held by the human;
the network remembers what the builder chooses to make public;
and the application can acquire users, history, reputation,
and economics without becoming dependent on a single company.
\end{proto}

\clearpage
% =====================================================================
% TOC
% =====================================================================
{\displayfont\LARGE Contents\par}
\vspace{1em}
\startcontents
\printcontents{}{1}{\setcounter{tocdepth}{1}}
\clearpage

% =====================================================================
\section{The name}
% =====================================================================

\textbf{\TOH{}} is \textbf{ONESHOT} reversed.

The reversal matters.

``One shot'' describes the act: a coherent intention is given a real
attempt at becoming software.

The reversal describes the thesis: the conventional software process is
turned around. Traditional software begins with infrastructure and
eventually tries to recover the original human intention from
requirements, tickets, meetings, and accumulated code. \TOH{} begins with
the intention and grows only the infrastructure that reality proves
necessary.

The name also encodes a discipline:

\begin{creed}
Do not confuse caution with fragmentation.\\
When the ontology is clear, carry it through the whole system.
\end{creed}

\TOH{} is not a promise that every application will be perfect after one
prompt. It is a promise that every coherent idea can receive an executable
first form, then continue evolving without losing its identity.

And beneath the thesis sits the invitation the whole system is built to
make easy, printed on the factory door and repeated at the end of every
successful build:

\begin{creed}
Take another one.
\end{creed}

% =====================================================================
\section{The problem}
% =====================================================================

\subsection{Human intentions are abundant; software projects are scarce}

People continuously imagine software that would improve their own lives: a
tool for a family ritual; a game for a niche community; a strange personal
instrument; a local business workflow; a training environment for a
specific skill; an interface that reflects one person's unusual way of
seeing the world.

Most of these ideas never become software. Not because they lack value,
but because the existing process assumes that software must justify the
cost of becoming a project.

\TOH{} begins from the opposite assumption:

\begin{creed}
Personal specificity is not a defect.\\
Weirdness is sufficient reason to create.
\end{creed}

\subsection{Creation is becoming easier than continuity}

Coding agents and increasingly capable development systems reduce the cost
of producing a working application. That is necessary, but insufficient.

A generated application still faces harder questions. Who owns it? How
does it remain understandable six months later? How does the builder
experience it on a real device? How does feedback return to the place
where the app evolves? How does one coherent app survive many subsequent
instructions? How can another person test it before it reaches the App
Store? How does the builder establish authorship? How do users carry
identity across apps without being universally tracked? How can the app
attach an economic object without binding itself to one launcher? What
remains if the company operating the tool disappears?

A generator answers, ``How do I produce code?''

\TOH{} asks, ``How does an intention acquire a life?''

\subsection{The precious-bet trap}

There is a subtler failure mode than not building at all: building one
thing and protecting it. When creating software is expensive, every idea
must be treated as the one precious bet. It gets a pitch deck before it
gets a body. It is defended before it is used. The builder learns to argue
for the idea instead of learning from it.

\TOH{} takes the photographer's position instead. Nothing has to be the
one precious bet. The unit of progress is the roll, not the frame. A Shot
is built to be used --- and allowed to die. The contact sheet, not the
pedestal.

\subsection{Existing platforms capture the application}

Many modern development platforms make creation easy by making the
resulting software dependent on their runtime, account system, hosting,
database, API keys, billing, or proprietary cloud. That convenience can
turn into captivity.

\TOH{} takes the opposite position:

\begin{creed}
The generated application belongs to the builder\\
before it belongs to the network.
\end{creed}

The network is additive. It provides continuity, not permission.

\subsection{Identity is either fragmented or invasive}

Most apps create their own accounts. The same human becomes a separate
database row everywhere. The apparent alternative is a universal
identifier exposed to every app, which enables effortless correlation and
surveillance.

\TOH{} seeks a third model: one underlying identity per human; a different
identity presented to each app; explicit, human-authorized proofs when
identities should be linked. Continuity should not require automatic
visibility.

\subsection{Distribution and economics are disconnected from creation}

An app can be technically complete and socially dead. App Store
distribution is powerful but slow, selective, and downstream of
substantial work. Open-source publication allows testing but lacks a
coherent discovery and reputation layer. Tokens can create attention and
economic motion, but token launch infrastructure is typically detached
from the software artifact it claims to represent.

\TOH{} connects these without making any one of them mandatory.

% =====================================================================
\section{The foundational primitive: the Shot}
% =====================================================================

\subsection{Definition}

A \textbf{Shot} is one coherent software intention embodied as an
independent application and preserved through its continuing lineage.

\begin{proto}
"Create a decision-training game for crypto traders."
-> one Shot

"Create an app that asks me to write before I scroll."
-> one Shot

"Create a tiny family app that remembers our bedtime stories."
-> one Shot
\end{proto}

A Shot is not: a Git commit; a prompt; a version; a workspace containing
several apps; a ticket; a release; a token; a child object beneath a
larger ``Project.''

Within the canonical \TOH{} ontology, the Shot is the application-level
object.

\subsection{The contact sheet}

The word \emph{Shot} is not a metaphor bolted on afterward. It is the
ontology.

In photography, a shot is one exposure: a deliberate act of attention,
committed to film, developed, and laid beside its siblings on a contact
sheet. The photographer does not publish the contact sheet. They study it.
Some frames get printed. Some teach the eye where to aim next. None of
them are failures; all of them are evidence.

A builder's collection of Shots is their contact sheet: a growing record
of intentions given bodies. The useful question about an idea is often not
``is this worth months?'' but ``what does it feel like when it exists?'' A
Shot makes that question cheap to answer. Held next to its siblings, it
does something a lone precious project never can:

\begin{creed}
Volume reveals signal.\\
Taste emerges through making.
\end{creed}

Some shots live. Some teach you where to aim next. The roll continues.

\subsection{One Shot, many evolutions}

The first coherent intent creates the Shot. Later instructions may
substantially transform it:

\begin{proto}
The Trenches
|-- first playable decision
|-- daily game
|-- post-game analysis
|-- rank
|-- community submissions
`-- App Store release
\end{proto}

These are Evolutions of one Shot. They are not six Shots. The original
circle remains around the same living object.

\subsection{Stable identity}

A Shot receives an immutable \texttt{shotId} when it is created.

The \texttt{shotId} must not be derived from mutable properties such as
the app name, the current source tree, the latest intent, the bundle
identifier, or the current owner. Names can change. Code can change.
Ownership can change. The Shot remains the same.

A Shot ID should therefore be generated once using sufficient entropy or
another collision-resistant method and retained in the Shot's local
manifest forever.

\subsection{Independence}

Every Shot must be independently operable.

A generated application must not secretly require: the \TOH{} CLI at
runtime; a \TOH{} account; an official \TOH{} API; the \TOH{} mobile app;
\texttt{\$TOHSENO}; a registry contract; or a proprietary build service.

A builder should be able to remove the global \TOH{} installation and
continue working on the Shot as an ordinary source repository.

This is not an implementation preference. It is constitutional.

\subsection{Ejectable from birth}

Independence has a stronger form, and the factory holds itself to it:
\textbf{ejection}. Ejection means the owner can build, operate, modify,
migrate, or replace a Shot without \TOH{}'s permission or a hidden factory
dependency --- and it is true at creation time, not a later export
feature.

Open source alone is insufficient. A repository can be public and still
captive. Ejection is made real by concrete properties every Shot carries
from its first commit: reproducible source; its own ordinary Git history;
pinned, locally runnable checks; explicit data and authority declarations;
owner-held identifiers; and no symlinks into a checkout, cache,
installation, or network package.

A Shot is \emph{not} ejectable if its core action requires a \TOH{}
account, endpoint, secret, CLI, mutable global validator, factory cache,
or undisclosed service; if build assets are linked out of the repository;
if only \TOH{} can read canonical data; or if leaving requires publishing
private input.

The factory's promise is symmetric: every Shot can leave, and the builder
can always take another one.

% =====================================================================
\section{Intent and Evolution}
% =====================================================================

\subsection{Intent}

An \textbf{Intent} is a human-readable statement of desired change.

The first intent answers: \emph{What should exist?}

A later intent answers: \emph{Given what reality taught me, what should
this Shot become next?}

An intent may be created from a direct instruction; feedback from a user;
observation of the app in use; a failed verification; a new design
decision; a change in the builder's taste; a market signal; a
conversation; or a constraint discovered on a real device.

The intent is private by default. It belongs to the builder's workshop,
not to the network, not to a provider's logs, and not to the generated
repository's public history. What crosses any boundary is a deliberate,
sanitized projection --- never the raw thought.

\subsection{Evolution}

An \textbf{Evolution} is a meaningful application of intent to an existing
Shot.

Not every file change is an Evolution. A formatter pass, dependency
update, or incidental commit may be part of an Evolution without
constituting one by itself. An Evolution marks a meaningful step in the
app's life: a change the builder would recognize as an answer to a real
intention. An Evolution is counted only when two independent judgments
agree: the coding agent completed its work, and deterministic verification
passed. Conceptually, an Evolution may contain:

\begin{proto}
shotId
evolution number
intent or private intent commitment
previous evolution hash
source tree hash
manifest hash
composition lock hash
artifact hash, when applicable
verification result
created time
builder signature, when made public
\end{proto}

The local repository may use Git as its detailed technical history.
\TOH{}'s Evolution history is the semantic history layered above it.

\subsection{The making loop}

\TOH{} is not complete when code appears. The full loop is:

\begin{proto}
intention
    |
    v
Shot
    |
    v
use on a real device
    |
    v
felt experience
    |
    v
feedback
    |
    v
sharpened next intent
    |
    v
Evolution
    |
    v
use again
\end{proto}

The purpose of continuity is to keep this loop intact. The app becomes a
conversation between intention and reality.

A shot becomes useful through contact with the real app, not through
terminal output alone. Ideas become true in use. And the loop has one more
arm, easily forgotten and deliberately honored: \emph{keep, evolve, or
kill it --- then take another shot.} Ending a Shot well is a legitimate
exit from the loop. Built to be used. Allowed to die.

% =====================================================================
\section{Lifecycle}
% =====================================================================

\subsection{Canonical states}

Every Shot has one canonical distribution state:

\begin{proto}
EVOLVING -> PUBLISHED -> APP_STORE
\end{proto}

The state is monotonic in public history. A later record may deprecate
availability, transfer control, or mark a release obsolete, but it must
not erase the fact that a Shot was previously published or released.

\subsection{EVOLVING}

A Shot is \textbf{EVOLVING} while it is primarily local to its builder.

It may be fully functional. It may run on a simulator or physical device.
It may have many Evolutions. It may even be shared privately. EVOLVING
does not mean unfinished in a pejorative sense. It means the Shot remains
inside the builder's workshop.

Properties: private by default; no network identity required; no public
source required; no blockchain required; continued evolution expected.

\subsection{PUBLISHED}

A Shot is \textbf{PUBLISHED} when a verifiable source representation is
made available through the \TOH{} ecosystem so that other people can
obtain and run it.

Published source should include enough information to answer: Which Shot
is this? Which Evolution is this? Where is the source? Which source tree
was published? How can it be built? Under what license is it available?
Which builder signed or claimed the publication?

Publication does not mean App Store release. It creates a space between
private development and centralized distribution where other people can
inspect the source, build the app, run it on their own device, test it,
provide feedback, fork it under its license, and verify its lineage.

A node may index and mirror a published Shot, but the node is not the
publisher. The builder publishes. The node witnesses, serves, and relays.

\subsection{APP\_STORE}

A Shot enters \textbf{APP\_STORE} when it is released through Apple's App
Store. In human language, this is the Shot being \textbf{released}.

An App Store record may include:

\begin{proto}
Apple app identifier
bundle identifier
version
source evolution
release artifact hash, when available
store URL
verification method
release time
\end{proto}

APP\_STORE is not the end state of creativity. A Shot may continue to
receive Evolutions forever while retaining APP\_STORE as its highest
distribution milestone.

The road from a built Shot to a live listing is walked by a human: an
Apple developer account, signing, a listing, review, and one step that is
genuinely irreversible --- \emph{submit for review; this one you cannot
take back}. \TOH{} may describe that ladder precisely, detect progress
honestly, and wait patiently. It must never pretend the ladder is shorter
than it is, and it must make irreversible rungs quieter, slower, and
clearer than ordinary steps.

\subsection{Lifecycle and recording are different}

A Shot's lifecycle state and its network recording status are orthogonal.
An EVOLVING Shot may record a private hash commitment proving that it
existed at a certain time. A PUBLISHED Shot may remain entirely off-chain.
An APP\_STORE Shot may use the network only for feedback. A Shot may have
a public builder record without an appcoin, or link an appcoin after
years of existence.

This distinction prevents the registry from becoming a gatekeeper.

% =====================================================================
\section{The local factory}
% =====================================================================

\subsection{The factory is the base layer}

\TOH{} must begin as a complete local tool. The expected experience is
conceptually simple:

\begin{proto}
install tohseno
tohseno            # take the first shot
tohseno create
tohseno evolve
tohseno verify
tohseno run
tohseno publish
\end{proto}

The exact command surface may evolve, but the principle does not:

\begin{creed}
A person can create and carry a Shot through its full software
lifecycle without joining the \TOH{} network.
\end{creed}

Speed is the product. The mission is intention to phone, measured in
minutes. Anything that adds a question, a config step, or a ceremony must
pay for itself in reliability.

\subsection{One intention, explicit interpretation}

The factory divides responsibility along a single, strict boundary:

\begin{creed}
Artificial intelligence interprets meaning.\\
Deterministic code owns composition, locks, storage invariants,\\
verification, and authority.
\end{creed}

The builder gives one private intention. A coding agent --- owner-selected
and replaceable --- reads it once, in a sandbox with no approvals and no
network write authority, and returns a small, sanitized, validated plan: a
name, a starting shape, a set of installed capabilities, a data posture,
an identity posture, a first definition of done. The plan is shown to the
human for approval before anything is built. If the agent is offline,
slow, or wrong, the factory falls back to a blank starting shape --- never
to another provider chosen on the builder's behalf.

From the approved plan onward, determinism takes over. Composition is a
locked overlay of a neutral kernel, a starting template, and
dependency-ordered capabilities, fingerprinted file by file. The
repository is materialized atomically: manifest, composition lock, pinned
verifier, provenance, Git baseline --- all created and verified before the
directory becomes visible. The agent then works \emph{inside} that
already-true repository, with a sanitized environment and bounded
authority, and its result is verified again after it exits. A failing
result is isolated, not presented as ready.

Capabilities are installed, not implied. Prose can guide an agent, but
prose cannot silently claim an uninstalled capability. If a feature cannot
be expressed as a valid manifest field, it is unsupported --- the factory
says so instead of improvising.

\subsection{Agent agnosticism}

\TOH{} may coordinate coding agents, models, compilers, templates, and
verification tools, but it must not define the identity of a Shot by which
model produced it.

A Shot can outlive the model used to create it, the agent harness, the
template version, the \TOH{} release, and the company operating any of
those systems. The protocol should record relevant composition
fingerprints without making a temporary intelligence provider the center
of the ontology.

\subsection{Truthful composition}

A Shot should carry enough local metadata to reconstruct how its current
form came to exist: app manifest; \TOH{} protocol version; template
identifiers; skill descriptors; lockfiles; immutable file hashes; verifier
results; source tree hashes; optional private provenance.

Private provenance must remain private by default. The public network
record is a deliberate projection of local truth, not a copy of the
entire workshop.

\subsection{Verification before ceremony}

\TOH{} should prefer executable checks over confidence language.

A Shot that claims to build should build. A Shot that claims to be
published should expose resolvable source. A Shot that claims an App Store
release should point to a verifiable listing. A Shot that claims an
appcoin should identify the chain, address, and deployment transaction.

The system should present verification levels honestly:

\begin{proto}
creator claimed
signature verified
source fingerprint verified
source available
release discovered
release control verified
token contract verified
deployment transaction verified
\end{proto}

Not all claims have equal strength. \TOH{} must never blur that
distinction.

The same honesty governs the moment of handoff. When a Shot is born, the
factory --- not the agent --- delivers the report: what was verified, what
was built, what launched, where the screenshot landed. The agent builds,
but \TOH{} is the only party allowed to say it worked, because its
statement is backed by executable evidence rather than enthusiasm. The
report ends with exactly one next action, and the reason for it:

\begin{proto}
NEXT
Use the app in Simulator and decide what its
next Evolution needs.

    tohseno my-app

WHY
A shot becomes useful through contact with the real app,
not through terminal output alone.
\end{proto}

\subsection{Open underneath, magic at the surface}

The factory's interior is radically inspectable: open source, pinned
versions, content-addressed releases, checksummed installs that re-verify
themselves on every invocation, and a repository gate that mechanically
enforces the system's own promises. Its surface is one line, one
intention, one working app.

These are not in tension. The magic is trustworthy precisely because the
underneath is boring, deterministic, and public. Transparent underneath.
Frictionless at the surface.

% =====================================================================
\section{How this document speaks: the honesty vocabulary}
% =====================================================================

A whitepaper written before a stable release faces a temptation: to
describe the dreamed system and the running system in the same tense.
\TOH{} refuses that tense. Throughout the project --- and throughout this
document --- claims are classified with a four-word vocabulary:

\begin{description}[leftmargin=1.6em, style=nextline]
\item[\statustag{IMPLEMENTED}] Exercised by a live code path and testable
now.
\item[\statustag{PREPARED}] Configuration and instructions exist, but no
external action has occurred.
\item[\statustag{PROPOSED}] Architecture or product behavior that is not
implemented.
\item[\statustag{OPEN}] Requires a product, security, or ownership
decision.
\end{description}

Applied to this document, at the moment of writing:

The Shot ontology, deterministic composition and locks, atomic
materialization, agent-boundary planning with blank fallback, pinned
verification, Evolution counting, Simulator run and capture, the Studio,
the checksummed installer, and protocol version 1's signed record schemas
with a reference node are \statustag{IMPLEMENTED}.

The human ladder from built Shot to App Store listing, and the
reference-node runbook, are \statustag{PREPARED}.

Publication records flowing from real builders, deployed nodes, persistent
builder identity, the mobile app, registry anchoring, and appcoin links in
the wild are \statustag{PROPOSED}.

The production trust root, cross-node fork resolution, anchoring
economics, the commercial boundary, and several questions collected in
\S\ref{sec:open} are \statustag{OPEN}.

Stale confident instructions are worse than an honest unknown. A future
reader comparing this document against a future system should be able to
tell exactly which promises were running, which were staged, and which
were still dreams --- and hold the project to the difference.

% =====================================================================
\section{Continuity}
% =====================================================================

\subsection{Definition}

\textbf{Continuity} is the ability of a Shot, a builder, and a user
relationship to survive changes in device, version, agent, node, host, and
distribution channel without surrendering sovereignty.

Continuity has three dimensions:

\begin{enumerate}
\item \textbf{Artifact continuity.} The Shot remains the same living
object across Evolutions and releases.
\item \textbf{Human continuity.} A person can carry identity, consent, and
selected relationships across apps and devices.
\item \textbf{Network continuity.} Public records and published artifacts
can remain available beyond one official server.
\end{enumerate}

\subsection{A continuity app}

{\sloppy
A \textbf{continuity app} is an independent application that may
optionally participate in \TOH{}'s continuity protocols. Participation may
allow an app to request: an app-specific user identity; a signed login
challenge; explicit cross-app linkage; feedback delivery; entitlement
proofs; recovery; a public Shot record; appcoin information; connection to
a builder's next-intent loop.\par}

An app remains functional if the user declines all of these.

The continuity layer is a capability, not a tax.

\subsection{The continuity paradox}

\TOH{} seeks two properties that are often treated as opposites:

\begin{proto}
the same human across many apps
        and
no automatic universal identifier across those apps
\end{proto}

The system resolves this through derived, app-specific identities and
explicit proofs of sameness.

Continuity without surveillance is not an optional privacy feature. It is
the identity model.

% =====================================================================
\section{Identity}
% =====================================================================

\subsection{There is no global seed phrase}

There is no global \TOH{} seed phrase shared by builders or users. Every
human controls an independent root.

\begin{proto}
JP root      =/=  Maria root
Alice root   =/=  Bob root
\end{proto}

No company-wide secret, protocol-wide mnemonic, or shared developer key
may stand in for personal identity.

\subsection{The Continuity Root}

A person may create a \textbf{Continuity Root}: a private cryptographic
root from which scoped identities can be derived. The root may be held by
the future \TOH{} mobile app; a local desktop signer; a hardware wallet; a
compatible third-party signer; or another open-source implementation.

The \TOH{} mobile app is a preferred premium custodian, not a protocol
monopoly. The root secret must never be transmitted to a node or generated
app.

\subsection{Builder identity}

A builder may derive or designate a stable public identity for authorship.
This identity can sign claims such as:

\begin{proto}
I created this Shot.
I authorize this Evolution as canonical.
I publish this source fingerprint.
I acknowledge this App Store release.
I link this appcoin to this Shot.
I transfer current control of this Shot.
\end{proto}

Builder identity is designed to accumulate public reputation. A builder
record may eventually show Shots created, Shots published, App Store
releases, verified source histories, community feedback, appcoin
relationships, node attestations. The purpose is not gamification for its
own sake. It is a credible body of work --- a public contact sheet.

\subsection{User identity}

When a person uses a continuity app, the signer derives an app-specific
identity. Conceptually:

\begin{proto}
APP_KEY = derive(
  CONTINUITY_ROOT,
  protocol domain,
  app identifier,
  environment
)
\end{proto}

The same person receives different identities in different apps:

\begin{proto}
Alice root
|-- Alice-in-Anky
|-- Alice-in-The-Trenches
`-- Alice-in-Garden-Memory
\end{proto}

Anky cannot automatically infer that \texttt{Alice-in-Anky} and
\texttt{Alice-in-The-Trenches} are the same person. This prevents the
continuity root from becoming a universal tracking number.

\subsection{Explicit cross-app proofs}

A person may choose to prove that two app-specific identities belong to
the same root. The signer presents a human-readable request:

\begin{proto}
THE TRENCHES wants to confirm
that you also use ANKY.

This links your identities between these two apps.

No other app identities will be revealed.
\end{proto}

Only after explicit approval does the signer produce a scoped continuity
proof. The proof should be specific to the requesting apps; specific to
the requested purpose; time-bounded where appropriate; revocable where the
claim is not permanently public; and incapable of revealing the root
secret.

\subsection{Builder and user roles}

The same human may be both a builder and a user. Those roles remain
distinct:

\begin{proto}
JP root
|-- public builder identity -> created Anky
`-- private app identity    -> JP-in-Anky
\end{proto}

An app does not automatically receive the builder's public identity merely
because the builder uses it.

\subsection{Collaboration}

Authorship, contribution, and control are different concepts. A Shot may
eventually support one genesis builder; several co-creators; credited
contributors; a current controller; an organization or multisignature
controller; a publisher-of-record for an App Store account.

Original authorship is historical and must not be silently transferred.
Current control may change. A future owner can control releases without
becoming the person who originally created the Shot.

\subsection{Recovery and rotation}

Continuity requires recovery, but recovery must not create an invisible
custodial backdoor. Compatible recovery methods may include encrypted
recovery packages; backups secured by the platform; user-chosen recovery
contacts; social recovery; key rotation signed by the current key; and
recovery policies held in contract accounts.

A rotated key may assume current control while old authorship signatures
remain valid historical facts.

% =====================================================================
\section{The \TOH{} mobile app and the Genesis invariant}
% =====================================================================

\subsection{What it is}

The future \TOH{} mobile app is an optional premium continuity client. It
may become: a continuity identity vault; a hardware-like signer; a mobile
approval surface; a portfolio of the builder's Shots; a place to
experience connected apps as a user; a structured feedback inbox; a
next-intent sharpening instrument; a bridge between real device use and
the local Studio; a hosted-network convenience layer; a subscription
product.

It is not the place where the source code fundamentally belongs. The
development machine remains the factory.

\subsection{What it is not}

The mobile app is not: required to install \TOH{}; required to create,
evolve, or publish a Shot; required to ship to the App Store; required for
a generated app to function; the only valid signer; the root owner of
generated apps; a mandatory wallet; or the protocol itself.

A builder must be able to use the complete open-source factory without
ever seeing the \TOH{} mobile app.

\subsection{Mobile authorization}

When paired with a local Studio, the phone may receive a readable request:

\begin{proto}
RECORD THIS SHOT

THE TRENCHES
Evolution 7

This will permanently publish:
  + Shot identifier
  + Builder identity
  + Source fingerprint
  + Lifecycle state

This will not publish:
  x Original private intent
  x Source files not already published
  x Agent conversation
  x Local secrets
\end{proto}

The builder may authorize with a deliberate gesture such as swiping right.
The gesture is not decoration. It represents:

\begin{creed}
I have seen the claim.\\
I understand its permanence.\\
Make it mine in public memory.
\end{creed}

\subsection{Feedback and the next intent}

The mobile app's most important product role may be the feedback loop. A
builder can use a Shot on the actual device where it matters, then
capture what felt wrong; what surprised them; what the app appears to want
to become; what another tester said; what should change next.

The mobile app can help compress this experience into a precise next
intent and send it to the local Studio. This is premium continuity:

\begin{proto}
the app remains close enough to life
that its next evolution is informed by use,
not by memory of what the builder thought before using it.
\end{proto}

\subsection{Subscription boundary}

\TOH{} may charge monthly or yearly for the mobile experience and hosted
convenience: synchronized builder identity; secure mobile signing;
cross-device recovery; hosted public builder profile; feedback routing;
intent sharpening; community discovery; node relay and gas sponsorship;
premium network indexing; continuity across connected apps.

It must not purchase the right to own, build, or continue operating one's
own source code.

\begin{creed}
\TOH{} charges for continuity, not freedom.
\end{creed}

\subsection{The genesis invariant}
\label{sec:genesis}

The \TOH{} mobile app must not be implemented by hand before the first
stable release of the protocol. The first stable \TOH{} system must create
it through the same public factory offered to every builder.

Genesis is generated. The mobile application's absence from the factory's
own repository is deliberate, and that absence is a tested property.
Building the app by hand, checking its source into the factory first, or
adding a privileged template that bypasses the ordinary path would
invalidate the proof that the factory can create its own product. The
mobile Shot's repository is an \emph{output} of the stable release, not an
input to it.

The intended sequence is:

\begin{proto}
stabilize the TOHSENO factory and protocol
        |
        v
run the stable `tohseno create`
        |
        v
create the TOHSENO mobile app as the Genesis Shot
        |
        v
evolve it through the ordinary lifecycle
        |
        v
publish its source
        |
        v
record its lineage
        |
        v
link its canonical appcoin
        |
        v
release it on the App Store
\end{proto}

No hidden template may special-case the Genesis Shot.

\begin{creed}
The loom must be complete before it weaves its first flag.
\end{creed}

% =====================================================================
\section{The network}
% =====================================================================

\subsection{The network is optional}

The network surrounds the local factory but does not replace it.

\begin{proto}
LOCAL FACTORY                     OPTIONAL NETWORK
-------------                     ----------------
create                            record
evolve                            mirror
verify                            discover
run                               route feedback
own                               relay
publish source                    link appcoin
ship                              preserve public memory
\end{proto}

If every \TOH{} node disappears, a builder's local Shots must remain
usable. Nodes index; builders own.

\subsection{The node}

A \textbf{Node} is a compatible service that performs one or more network
functions. A node may: receive signed public records; validate schemas;
verify signatures and hashes; store or pin public manifests; mirror
published source; index registry events; serve builder and Shot profiles;
route encrypted feedback; relay transactions; sponsor gas; verify appcoin
deployments; verify public release claims; synchronize public records with
peers.

A node does not become the author, owner, or publisher merely by serving a
Shot. A node is an index and a transport, not an ownership authority or a
consensus system.

\begin{proto}
The builder authors.
The builder publishes.
The node witnesses and serves.
The registry remembers.
\end{proto}

\subsection{The official node}

The first official \TOH{} service may run a minimal server. It can provide
the default experience, but all clients should depend on a documented node
interface rather than hardcoding one domain --- a rule the factory
enforces mechanically against its own code. Conceptually:

\begin{proto}
interface TohsenoNode {
  submitRecord(record: SignedRecord): Promise<RecordReceipt>;
  getShot(shotId: ShotId): Promise<NetworkShot | null>;
  getBuilder(id: PublicIdentity): Promise<BuilderRecord | null>;
  linkAppcoin(link: SignedAppcoinLink): Promise<LinkReceipt>;
  routeFeedback(payload: EncryptedFeedback): Promise<DeliveryReceipt>;
}
\end{proto}

A user should eventually be able to choose:

\begin{proto}
tohseno network use https://node.tohseno.com
tohseno network use https://community.example
tohseno network use http://localhost:7036
\end{proto}

The abstraction is small by design. It leaves the door open to
decentralization without pretending decentralization already exists.

\subsection{Objective verification, not aesthetic judgment}

Nodes may verify objective properties: the signature matches the claimed
builder; the manifest hashes to the claimed value; the source tree matches
the public source; the parent record exists; the token exists at the
claimed address; the deployment transaction created the claimed token; the
App Store identifier resolves; the record follows the protocol schema.

Nodes should not initially decide whether an app is good; whether an idea
is original enough; whether the creator ``deserves'' authorship; whether a
token is economically wise; whether an app is culturally important.

\TOH{} records claims and evidence. It does not claim omniscience.

\subsection{Existing consensus, not a new blockchain}

\TOH{} does not need to invent its own blockchain to decentralize nodes.
An existing chain can provide ordering, timestamps, transaction finality,
and globally readable events. Nodes provide the surrounding work: storage,
indexing, verification, availability, relaying, discovery.

This keeps the protocol narrow.

\subsection{Progressive decentralization}

The likely path is:

\begin{description}[leftmargin=1.6em, style=nextline]
\item[Stage 0 --- Local-only.] The factory works with no node.
\item[Stage 1 --- Official node.] Anky, Inc.\ or another initial steward
operates the first node.
\item[Stage 2 --- Community mirrors.] Anyone can run a node that syncs
public records, mirrors manifests, and serves discovery.
\item[Stage 3 --- Active relayers.] Multiple authorized or bonded nodes
can validate and relay new records.
\item[Stage 4 --- Economic coordination.] Nodes may stake and earn
\texttt{\$TOHSENO} for measurable, verifiable work.
\end{description}

This sequence matters. A token must not subsidize empty infrastructure in
the hope that meaningful applications eventually appear.

\begin{creed}
The organism must breathe before its metabolism is financialized.
\end{creed}

% =====================================================================
\section{Public memory and the registry}
% =====================================================================

\subsection{The purpose of the registry}

The registry answers small, durable questions:

\begin{proto}
Did this builder claim this Shot?
Did this Evolution belong to this Shot?
Did the builder publish this source fingerprint?
Did this Shot reach a lifecycle milestone?
Which appcoin did the builder declare canonical?
Who currently controls official updates?
What was the public history of those claims?
\end{proto}

It does not store the entire application.

\subsection{Append-only facts}

The registry is append-only. A record may supersede an earlier record, but
must not erase it. This permits key rotation, appcoin migration, control
transfer, deprecation, corrected metadata, and new lifecycle milestones.
History remains visible.

Records are portable builder attestations, not network truth. An honest
registry knows the difference and says so.

\subsection{Record types}

A minimal protocol may define records such as:

\begin{proto}
SHOT_CLAIM
EVOLUTION_CLAIM
LIFECYCLE_TRANSITION
APPCOIN_LINK
CONTROL_TRANSFER
DEPRECATION
\end{proto}

The exact encoding may evolve through protocol versions, but the semantic
boundaries should remain clear. Unknown fields fail validation. There are
deliberately no fields for raw prompts, private provenance digests, local
databases, app-user content, conversations, credentials, secrets,
unpublished source bytes, or model reasoning. A safe-looking summary can
still disclose an idea, so signing and submitting a record must remain a
deliberate public action.

\subsection{A generic record}

Conceptually:

\begin{proto}
{
  "protocol": "tohseno",
  "protocolVersion": 1,
  "recordId": "0x...",
  "recordType": "LIFECYCLE_TRANSITION",
  "shotId": "0x...",
  "previousRecordId": "0x...",
  "actor": "0xBuilder...",
  "manifestHash": "0x...",
  "manifestURI": "ar://...",
  "createdAt": "2026-07-26T16:00:00Z"
}
\end{proto}

The \texttt{recordId} should commit to the exact canonical representation.
The human-readable manifest may live in a \TOH{} node, object storage,
IPFS, Arweave, or another content-addressed system. The onchain record
stores a hash and a pointer.

\subsection{Creator authorization}

A public claim should be authorized by the builder or current controller.
A typed-data signature can bind: protocol domain; registry; chain; record
type; Shot ID; manifest hash; nonce; expiration. The signature should be
human-readable before approval. A node may submit and pay gas without
becoming the creator.

\subsection{Gasless experience}

The preferred experience is:

\begin{proto}
review -> sign -> recorded
\end{proto}

Not:

\begin{proto}
buy ETH -> bridge -> select network -> estimate gas -> retry
\end{proto}

A node may sponsor gas for eligible records. A builder may also self-relay
or use another node. Gas sponsorship is a service, not a protocol
dependency.

\subsection{Why the registry is not an NFT}

The permanent authorship record should not begin as a transferable NFT.
Authorship, control, ownership, and collectibility are different concepts.
An ordinary transferable NFT can imply: \emph{whoever holds this token is
the creator or owner of the represented app.} That may be false.

The registry should preserve: \emph{this identity created or authorized
this public record at this time.} Future transferable project-control
objects or collectibles may exist on top of that history, but they must
not rewrite authorship.

\subsection{Public record as career memory}

The registry is useful because it can become the memory of a builder's
life in software. A credible profile may eventually show:

\begin{proto}
the Shots a person created
how long they carried them
which ones reached publication
which ones reached the App Store
which appcoins were linked
which source histories were verified
which communities formed around them
\end{proto}

\begin{creed}
Provenance is not paperwork.\\
It is a body of work that can survive a platform.
\end{creed}

% =====================================================================
\section{Public and private data}
% =====================================================================

\subsection{Local truth, public projection}

The local Shot contains richer information than the public network record.
The public record must be derived intentionally. A publication preview
should state exactly what crosses the boundary.

\subsection{Public by explicit choice}

Possible public fields: Shot ID; app name; icon; builder public identity;
lifecycle state; Evolution number; source locator; source tree hash;
manifest hash; build instructions; license; App Store identifiers; appcoin
chain and address; public screenshots; public description.

\subsection{Private by default}

Private fields include: unpublished source; original private intentions;
agent conversations; local databases; secrets; signing keys; user data;
private feedback; internal analytics; device identifiers; recovery
material.

Private-by-default and account-free are defaults, not refusals. A manifest
may truthfully declare other data, identity, integration, entitlement, or
irreversible mechanics. Those are product decisions the builder makes in
the open --- streaks, paywalls, and scores are tools, not sins --- and the
manifest exists so that such declarations are explicit rather than
implied.

\subsection{Hashes are not automatically private}

A hash of low-entropy text can leak information through guessing. Private
intentions should use appropriate commitments, including salts or keyed
constructions, when later proof is desired. \TOH{} must not call data
``private'' merely because it has been hashed.

\subsection{Irreversibility requires legibility}

Before any permanent record, the user must be able to understand: what
will become public; what will remain private; which identity signs;
whether a node will sponsor gas; whether the record can be superseded;
whether the action can ever be forgotten.

The protocol should make irreversible actions quieter, slower, and clearer
than ordinary interface taps.

% =====================================================================
\section{Appcoins}
% =====================================================================

\subsection{Definition}

An \textbf{Appcoin} is an optional economic artifact canonically linked to
a Shot. It may represent a community's shared attention; a way to fund
continued development; an instrument for distribution; access, membership,
or coordination; a market around an app's cultural motion; an experiment
whose meaning evolves with the Shot.

An appcoin is not the Shot. The Shot can exist without a coin. The coin
should point back to a Shot.

Recording a link does not deploy, mint, trade, transfer, price, or endorse
an asset. \TOH{} never promises wealth, and never makes financial
mechanics, urgency, or speculation the reason to build. The prototype is
the payoff; the coin, when it exists, is downstream of a living thing.

\subsection{Deployment agnosticism}

\TOH{} must not require appcoins to be deployed by the registry contract.
A coin may be deployed through any launcher, a custom smart contract, a
future chain, or a mechanism not yet invented. The \TOH{} record links the
canonical relationship:

\begin{proto}
Shot ID
coin chain ID
token address
deployment transaction hash
deployment provider, if known
linking builder signature
link time
economic metadata
\end{proto}

The protocol cares about the relationship, not the launcher.

\subsection{Canonical appcoin}

A Shot may have zero or one active canonical appcoin at a time. If a coin
migrates, the new link explicitly supersedes the prior one; the old
relationship remains historical. Because Evolutions belong to the same
Shot, an appcoin does not need to be recreated for every version.

\subsection{Verification in both directions}

A strong appcoin relationship can point both ways:

\begin{proto}
TOHSENO record            -> token address
token deployment metadata -> TOHSENO Shot ID
\end{proto}

A one-way builder claim is still useful, but a bidirectional link is
stronger.

\subsection{Fee capture}

A deployment-agnostic registry cannot magically collect fees from
arbitrary token contracts. Any fee routing must be explicit. Possible
future models include a launch adapter that routes a declared portion of
trading fees; a token contract that names builder, launcher, and network
beneficiaries; a voluntary protocol fee; a node relay fee; a
subscription-funded public service; app-specific economic rules.

\TOH{} must never claim revenue rights that are not encoded and consented
to.

\subsection{\texttt{\$TOHSENO}}

The intended genesis is for \texttt{\$TOHSENO} to become the first
canonical appcoin recorded by the stable protocol and linked to the \TOH{}
Genesis Shot. Its deployment mechanism remains replaceable. Its potential
roles are distinct:

\begin{enumerate}
\item \textbf{Appcoin of the \TOH{} mobile app.} It is economically
connected to the first canonical app created by the system.
\item \textbf{Coordination asset for the network.} It may eventually be
staked or earned by nodes performing verifiable work.
\item \textbf{Shared economic surface of the ecosystem.} Future appcoin
adapters may choose to route explicit value toward the network or
\texttt{\$TOHSENO}.
\end{enumerate}

None of these roles may make \texttt{\$TOHSENO} mandatory for local
creation, ordinary app use, source ownership, or App Store release.

\subsection{Node rewards}

If node rewards are introduced, they should compensate measurable work:
storing and serving manifests; mirroring published source; validating
records; relaying transactions; indexing events; maintaining availability;
providing cryptographic proofs of service.

Rewards should not be issued merely because a process claims to be online.
Otherwise, the network invites empty node farming rather than useful
infrastructure. Tokenomics should describe real metabolism, not substitute
for it.

% =====================================================================
\section{The end-to-end system}
% =====================================================================

\subsection{Local creation \hfill \statustag{IMPLEMENTED}}

\begin{proto}
Builder writes first intent
        |
        v
agent interprets it into a sanitized, approvable plan
        |
        v
TOHSENO composes deterministically and creates a Shot
        |
        v
Shot receives immutable Shot ID
        |
        v
independent source repository, atomic from birth
        |
        v
local run, verification, and Simulator capture
\end{proto}

No account, node, mobile app, token, or chain is involved.

\subsection{Evolution \hfill \statustag{IMPLEMENTED}}

\begin{proto}
Builder uses the app
        |
        v
captures feedback
        |
        v
sharpens next intent
        |
        v
TOHSENO applies an Evolution
        |
        v
verification runs
        |
        v
same Shot, new form
\end{proto}

\subsection{Publication \hfill \statustag{PROPOSED}}

\begin{proto}
Builder chooses PUBLISHED
        |
        v
TOHSENO prepares public projection
        |
        v
source tree and manifest are fingerprinted
        |
        v
builder signs locally or with mobile signer
        |
        v
source is made available
        |
        v
any compatible node can index or mirror it
\end{proto}

The builder may publish source without an onchain record.

\subsection{Permanent recording \hfill \statustag{PROPOSED}}

\begin{proto}
factory prepares readable claim
        |
        v
signer shows exact public fields
        |
        v
human authorizes
        |
        v
node validates signature and schema
        |
        v
node stores manifest
        |
        v
node relays hash to registry
        |
        v
public receipt returns
\end{proto}

\subsection{App Store release \hfill \statustag{PREPARED}}

\begin{proto}
Shot reaches Apple release
        |
        v
release metadata is captured
        |
        v
builder signs APP_STORE transition
        |
        v
node verifies available evidence
        |
        v
registry records milestone
\end{proto}

The Shot continues evolving afterward.

\subsection{Appcoin link \hfill \statustag{PROPOSED}}

\begin{proto}
builder deploys coin anywhere
        |
        v
TOHSENO receives chain, address, deployment transaction
        |
        v
builder reviews canonical link
        |
        v
builder signs
        |
        v
node verifies token and transaction
        |
        v
registry appends APPCOIN_LINK
\end{proto}

\subsection{User continuity \hfill \statustag{PROPOSED}}

\begin{proto}
User opens a continuity app
        |
        v
app requests scoped identity
        |
        v
compatible signer derives app-specific key
        |
        v
user approves
        |
        v
app receives proof, not root secret
\end{proto}

\subsection{Cross-app continuity \hfill \statustag{PROPOSED}}

\begin{proto}
App A requests proof of relationship to App B
        |
        v
signer shows exact scope
        |
        v
user approves
        |
        v
signer creates pair-specific proof
        |
        v
only the requested identities become linkable
\end{proto}

% =====================================================================
\section{A minimal protocol model}
% =====================================================================

The following model is illustrative, not a frozen wire format.

\subsection{Local Shot manifest}

\begin{proto}
{
  "schema": "tohseno.shot/v1",
  "shotId": "shot_01...",
  "name": "The Trenches",
  "createdAt": "2026-07-26T16:00:00Z",
  "lifecycle": "EVOLVING",
  "currentEvolution": 7,
  "platform": "ios",
  "bundleId": "app.thetrenches.mobile",
  "protocolVersion": "1.0.0",
  "composition": {
    "template": "ios-native",
    "skillsLockHash": "0x...",
    "sourceTreeHash": "0x..."
  },
  "network": {
    "recorded": false,
    "canonicalNode": null
  }
}
\end{proto}

\subsection{Evolution manifest}

\begin{proto}
{
  "schema": "tohseno.evolution/v1",
  "shotId": "shot_01...",
  "evolution": 7,
  "previousEvolutionHash": "0x...",
  "intentCommitment": "0x...",
  "sourceTreeHash": "0x...",
  "manifestHash": "0x...",
  "artifactHash": null,
  "verification": {
    "status": "passed",
    "reportHash": "0x..."
  },
  "createdAt": "2026-07-26T16:30:00Z"
}
\end{proto}

\subsection{Public Shot record}

\begin{proto}
{
  "schema": "tohseno.public-record/v1",
  "recordType": "LIFECYCLE_TRANSITION",
  "shotId": "shot_01...",
  "evolution": 7,
  "builder": "0x...",
  "lifecycle": "PUBLISHED",
  "source": {
    "uri": "https://...",
    "treeHash": "0x...",
    "license": "..."
  },
  "manifestHash": "0x...",
  "previousRecordId": "0x...",
  "createdAt": "2026-07-26T16:45:00Z"
}
\end{proto}

\subsection{Appcoin link}

\begin{proto}
{
  "schema": "tohseno.appcoin-link/v1",
  "shotId": "shot_01...",
  "builder": "0x...",
  "chainId": 8453,
  "tokenAddress": "0x...",
  "deploymentTxHash": "0x...",
  "deploymentProvider": "bankr",
  "supersedes": null,
  "createdAt": "2026-07-26T17:00:00Z"
}
\end{proto}

\subsection{Signing request}

\begin{proto}
{
  "schema": "tohseno.signing-request/v1",
  "action": "LINK_APPCOIN",
  "humanSummary": {
    "title": "Link official appcoin",
    "shotName": "The Trenches",
    "publicFields": [
      "Shot identifier",
      "Builder identity",
      "Chain ID",
      "Token address",
      "Deployment transaction"
    ]
  },
  "payloadHash": "0x...",
  "expiresAt": "2026-07-26T17:10:00Z"
}
\end{proto}

\subsection{Node receipt}

\begin{proto}
{
  "schema": "tohseno.node-receipt/v1",
  "node": "node_public_key",
  "recordId": "0x...",
  "manifestHash": "0x...",
  "storage": {
    "uri": "ar://...",
    "status": "stored"
  },
  "registry": {
    "chainId": 8453,
    "transactionHash": "0x...",
    "confirmations": 12
  },
  "receivedAt": "2026-07-26T17:01:00Z"
}
\end{proto}

% =====================================================================
\section{Security model}
% =====================================================================

\subsection{Nodes are not trusted with identity secrets}

A node receives public keys, signed claims, public manifests, and
encrypted payloads. It must never receive the Continuity Root;
app-specific private keys; recovery phrases; unencrypted private feedback
unless explicitly addressed to the node as a service; or unpublished
source by default.

\subsection{Generated apps are not trusted with the root}

A generated app receives an app-specific proof or asks the signer to sign
a challenge. It should not receive the Continuity Root. A compromised
experimental app must not compromise every other continuity app.

\subsection{Agents are not trusted with the factory}

The factory extends the same discipline to the coding agents it
coordinates. An agent plans in a read-only sandbox and builds inside a
repository whose protected state is snapshotted before and verified after
its run. Its environment is sanitized: credentials the factory happens to
inherit are never forwarded merely because the agent would find them
convenient. An agent's claim of success is not accepted; it is checked. A
result that fails verification is isolated, not shipped.

\subsection{Domain separation}

Derived keys and signatures must use explicit domains. A login signature
must not be replayable as an appcoin authorization. A signature for one
registry or chain must not silently authorize another. A test environment
must not share production authority.

\subsection{Replay protection}

Permanent and temporary signing requests should include appropriate
nonces, deadlines, record IDs, previous-record references, and chain and
contract domains.

\subsection{Canonical encoding}

Hashing requires deterministic encoding. The protocol must define field
order; Unicode normalization; number representation; timestamp
representation; omission versus null; URI normalization; canonical JSON or
another deterministic format. Without canonicalization, two honest
implementations can disagree about the same record.

\subsection{Local-first does not mean careless}

Local keys require secure system storage; explicit export; backups;
rotation; clear environment separation; no plaintext secrets in
repositories. The mobile signer may improve key safety, but must not be
the only safe path.

\subsection{Supply chain honesty}

The path from ``curl'' to ``factory'' is part of the security model. An
installer should pin exact versions and checksums; refuse insecure
transport; refuse to mutate installations it cannot authenticate; and
leave behind a tool that re-verifies its own integrity every time it
runs. Magic at the surface is only tolerable when the surface can prove
what is underneath.

\subsection{Open-source verification}

The security of \TOH{} should not depend on users trusting a closed
official client. Protocol schemas, signature logic, registry contracts,
and reference node behavior should be inspectable.

\subsection{Registry humility}

An onchain timestamp proves that a signed claim existed by a certain
block. It does not automatically prove the builder was the first human
ever to imagine the idea; the source was free of copied work; the App
Store account was legally owned; the token was safe; or the app was
valuable.

\TOH{} must state exactly what each proof establishes and no more.

% =====================================================================
\section{Sustainability}
% =====================================================================

\subsection{What remains free}

The following must remain available without a subscription: installing
the open-source factory; creating Shots; evolving Shots; running and
verifying source; owning generated repositories; publishing source through
compatible mechanisms; using alternative signers; using alternative nodes;
self-hosting; shipping independently.

\subsection{What can be paid}

\TOH{} may charge for services that create real ongoing cost or
convenience: mobile continuity; hosted identities and recovery assistance;
synchronized feedback; sharpened next-intent workflows; official-node
storage; relaying and sponsored gas; public profiles; discovery; community
testing; managed publication; managed App Store operations; premium agent
orchestration; organization features.

The business should grow around the open protocol rather than weakening
it.

\subsection{Builder-aligned economics}

If \TOH{} participates economically in appcoins, the relationship must be
explicit. The builder should be able to see who receives fees; why; at
what rate; through which contract or provider; and whether the
relationship can change. Hidden extraction would contradict the protocol's
purpose.

\subsection{The network as a service before a market}

Before nodes receive token rewards, the official node should prove that
builders want public records; users want published apps; feedback actually
returns; storage and relay work are measurable; and independent nodes can
reproduce the behavior.

% =====================================================================
\section{Protocol evolution and governance}
% =====================================================================

\subsection{Versioned specifications}

\TOH{} protocol artifacts must be versioned:

\begin{proto}
shot schema
evolution schema
public record schema
signing protocol
node API
identity derivation
registry contract
continuity SDK
\end{proto}

Old records remain interpretable forever.

\subsection{Backward compatibility}

A new \TOH{} release should prefer additive fields; explicit migrations;
versioned endpoints; stable record semantics; preserved local operability.
A future agent must not silently reinterpret an old \texttt{PUBLISHED}
record to mean something new.

\subsection{Constitutional changes}

Some implementation details may change freely. The core invariants require
exceptional scrutiny. A protocol version that makes the mobile app
mandatory, turns the official node into a gatekeeper, exposes a universal
user identifier, or makes generated apps dependent on \TOH{} at runtime is
not a routine upgrade. It is a different system.

\subsection{Stewardship}

The initial repository and official infrastructure may have a central
steward. That stewardship is practical, not metaphysical. The long-term
credibility of \TOH{} comes from open specifications; independent
implementations; portable signatures; content-addressed records;
self-hostable nodes; sovereign generated apps.

\subsection{Forkability}

If a future steward violates the constitutional invariants, builders
should be able to fork the factory, node software, protocol schemas,
public indexes, registry readers, and continuity clients.

\begin{creed}
A protocol that cannot survive disagreement\\
is a product, not a protocol.
\end{creed}

% =====================================================================
\section{The Genesis Shot}
% =====================================================================

\subsection{The bootstrap problem}

The \TOH{} mobile app is intended to become the premium identity and
signing surface for future public records. But it cannot authorize its own
creation before it exists. The protocol therefore bootstraps honestly:

\begin{enumerate}
\item The pre-stable system uses a local or external signer.
\item The factory reaches its first stable release.
\item That stable release creates the \TOH{} mobile app through the
ordinary \texttt{tohseno create} path.
\item The app evolves through the ordinary Shot lifecycle.
\item Its source is published.
\item Its public lineage is signed using the bootstrap signer.
\item Once the mobile signer is operational, control can rotate or extend
to it.
\item \texttt{\$TOHSENO} is linked through the same appcoin mechanism
offered to future builders.
\item The app reaches the App Store through the same lifecycle semantics
used by every later Shot.
\end{enumerate}

The bootstrap is part of the history, not an embarrassment to hide.

\subsection{No special template}

The stable factory must be capable of creating the \TOH{} mobile app
without private generators or hand-written scaffolding unavailable to
other builders. It may use ordinary public skills, templates, and platform
capabilities. The Genesis Shot should demonstrate the system, not receive
an exemption from it.

Forbidden before genesis: a mobile app target in the factory repository; a
prebuilt mobile repository waiting in the wings; a privileged template
that bypasses the ordinary path; a mobile-only identity shortcut; a
fabricated genesis Shot ID or evidence record.

\subsection{Genesis readiness}

The first stable release is not ready merely because the CLI runs. It is
ready when it can create and evolve the class of application required by
the Genesis Shot, including native iOS foundations; secure local storage;
networking; pairing; signing abstractions; subscriptions; feedback
interfaces; Shot browsing; protocol integration; App Store readiness; and
independent source ownership.

\subsection{The recursive closure}

The intended recursion is:

\begin{proto}
TOHSENO creates the TOHSENO app
        |
        v
the TOHSENO app holds continuity identity
        |
        v
continuity identity authorizes future TOHSENO records
        |
        v
those records preserve the lineage of future Shots
        |
        v
some of those Shots create new continuity apps
\end{proto}

The system becomes capable of remembering the objects it helps bring into
existence.

% =====================================================================
\section{What \TOH{} is not, and the negative space}
% =====================================================================

\TOH{} is not a promise that every idea deserves a market. It is a promise
that a coherent idea deserves a real attempt at embodiment.

\TOH{} is not: a mandatory cloud runtime; a centralized app account; an
App Store replacement; a token launcher disguised as a developer tool; an
NFT minting platform; a universal user tracking system; a new blockchain;
a DAO before a product; a node-reward farm; an agent-specific wrapper; a
repository of everyone's private prompts; a system that owns generated
source; a social network that must be joined before creating; a claim that
code generation alone is continuity.

It may interoperate with many of these systems. It must not lose its
identity inside them.

There is a deeper pattern here worth naming. \TOH{} defines itself as much
by absence as by presence, and it treats absence as a property to be
enforced, not merely intended. There is no global seed phrase --- by rule.
There is no hand-built mobile app --- by rule, until genesis. There are no
protocol fields for private prompts --- by schema. There is no fabricated
usage number --- by brand contract. The official hostname may not be
hardcoded into replaceable protocol code --- by mechanical check. In the
places where most systems accumulate quiet dependencies and quiet
exaggerations, \TOH{} keeps tested emptiness.

\begin{creed}
Absence, here, is a tested property.
\end{creed}

% =====================================================================
\section{Constitutional invariants}
% =====================================================================

The following statements define \TOH{} more deeply than any
implementation.

\subsection{The Shot}
\begin{enumerate}[itemsep=0.1em]
\item One coherent intention creates one Shot.
\item A Shot is the application-level object.
\item A Shot evolves.
\item An Evolution is not a new Shot.
\item The Shot ID remains stable across names, versions, owners, and
releases.
\end{enumerate}

\subsection{Lifecycle}
\begin{enumerate}[itemsep=0.1em, resume]
\item The canonical lifecycle is \texttt{EVOLVING} $\rightarrow$
\texttt{PUBLISHED} $\rightarrow$ \texttt{APP\_STORE}.
\item PUBLISHED means verifiable source is available through the
ecosystem.
\item APP\_STORE means released through Apple.
\item A Shot may continue evolving after every lifecycle milestone.
\item Recording and lifecycle are orthogonal.
\end{enumerate}

\subsection{Sovereignty}
\begin{enumerate}[itemsep=0.1em, resume]
\item The local factory requires no account, mobile app, wallet, token,
official node, or blockchain.
\item Every generated Shot remains independently operable and ejectable
from birth.
\item \TOH{} network features are additive.
\item Builders own their source.
\item The official service must not become a kill switch over local apps.
\end{enumerate}

\subsection{Identity}
\begin{enumerate}[itemsep=0.1em, resume]
\item There is no global \TOH{} seed phrase.
\item Every human controls an independent root.
\item Generated apps never receive the root secret.
\item App-specific identities are distinct by default.
\item Cross-app linkage requires explicit human consent.
\item Builder identity and user identity are different roles.
\item Authorship history cannot be silently transferred.
\end{enumerate}

\subsection{Mobile}
\begin{enumerate}[itemsep=0.1em, resume]
\item The \TOH{} mobile app is optional.
\item It may act as a premium identity vault, signer, feedback layer, and
continuity client.
\item It must not be required for open-source creation or app operation.
\item It must be the Genesis Shot created by the first stable release.
\item It must not be hand-built into the pre-stable repository.
\end{enumerate}

\subsection{Network}
\begin{enumerate}[itemsep=0.1em, resume]
\item A node witnesses, verifies, stores, serves, and relays.
\item A node is not the builder or publisher.
\item The official node is the first node, not the network.
\item Clients should be able to use compatible alternative nodes.
\item Nodes verify objective evidence rather than artistic worth.
\item Existing chains provide consensus; \TOH{} need not invent one.
\end{enumerate}

\subsection{Registry}
\begin{enumerate}[itemsep=0.1em, resume]
\item Public records are append-only.
\item Records contain small signed facts and content hashes.
\item Private source and secrets do not belong onchain.
\item Nodes may sponsor gas without becoming authors.
\item Authorship records are not ordinary transferable NFTs.
\item Public claims must say exactly what they prove.
\end{enumerate}

\subsection{Appcoins}
\begin{enumerate}[itemsep=0.1em, resume]
\item Appcoins are optional.
\item An appcoin is linked to a Shot; it is not the Shot.
\item Deployment remains launcher-agnostic.
\item The registry does not need to deploy tokens.
\item Fee routing must be explicit.
\item \texttt{\$TOHSENO} must never be required to create, own, or run a
Shot.
\end{enumerate}

\subsection{Continuity and truth}
\begin{enumerate}[itemsep=0.1em, resume]
\item The full loop is intention $\rightarrow$ Shot $\rightarrow$ use
$\rightarrow$ feedback $\rightarrow$ next intent $\rightarrow$ Evolution.
\item Continuity should preserve the human relationship to the app, not
merely its files.
\item Artificial intelligence interprets meaning; deterministic code owns
composition, verification, and authority.
\item Claims are classified honestly: implemented, prepared, proposed, or
open --- and the system never blurs them.
\item The protocol must remain understandable by future humans and
machines; old public records remain interpretable; and the system must be
forkable if its steward abandons these invariants.
\end{enumerate}

% =====================================================================
\section{Open questions}
\label{sec:open}
% =====================================================================

This document fixes the ontology, not every implementation decision. The
project keeps its unresolved decisions in a public list whose stated goal
is its own deletion: \emph{the asymptote of that file is empty.} At the
time of writing, the following remain deliberately open:

\begin{itemize}
\item The path currently stops at Simulator: how should the factory
accompany the human ladder to a physical device and a live listing without
pretending to automate what it cannot?
\item Which canonical chain should host the first registry?
\item Which signature scheme best balances wallet compatibility, passkeys,
and device security?
\item What is the production Builder trust root, and how should competing
valid histories for the same Shot ID be resolved across nodes? (Protocol
version 1 has no fork-resolution rule.)
\item How should source availability be proven across independent nodes?
\item Which manifest storage should be default before decentralized
pinning is mature?
\item How should organization and co-creator control be represented?
\item What constitutes strong App Store ownership verification?
\item How should key recovery work without making the official mobile app
custodial?
\item Which feedback should be public, private, encrypted, or anonymous?
\item How should a published Shot declare licensing and fork
relationships?
\item How should appcoin migrations and compromised coins be handled?
\item Which measurable node services deserve rewards, and how does the
network prevent Sybil farming?
\item Should nodes stake \texttt{\$TOHSENO}, earn it, burn it, or use
another coordination model?
\item How can a user prove cross-app facts without revealing linkable
identifiers?
\item Which continuity capabilities belong in every generated app
template?
\item Can any commercial layer be added while preserving the account-free
local factory?
\item What exactly must the stable release demonstrate before the Genesis
Shot is created?
\end{itemize}

These questions should be answered by preserving the invariants, building
the smallest truthful implementation, and observing real use.

% =====================================================================
\section{The wager}
% =====================================================================

\TOH{} makes a wager about the future of software.

It wagers that there will not merely be more applications. There will be
more \textbf{personal applications}: software specific enough to embody
the taste, rituals, needs, obsessions, relationships, and strange internal
logic of individual people.

It wagers that a person may carry dozens or hundreds of living Shots ---
a contact sheet that grows for a lifetime.

It wagers that the cost of creating a first form will continue falling,
and that the scarce things will become: coherent intention; taste;
attention; honest feedback; continuity; trust; distribution; public
memory.

It wagers that most shots will miss, and that this is precisely why the
system works: because a missed shot that exists teaches more than a
perfect shot that was never taken, and because taste is developed on the
contact sheet, not in the pitch deck.

It wagers that builders will want the freedom of local ownership and the
benefits of a network without being forced to choose between them.

It wagers that identity can be continuous without becoming a surveillance
identifier.

It wagers that an app can acquire an economic life without its
token-launch mechanism becoming its ontology.

It wagers that the most important future software may begin not as a
startup, but as one person saying:

\begin{creed}
I wish this existed.
\end{creed}

\TOH{} gives that sentence a Shot.

% =====================================================================
\section{Closing}
% =====================================================================

A Shot begins as a circle drawn around an intention.

Inside the circle, the idea becomes executable. Then reality touches it.
The app is used. It resists. It surprises. It teaches the builder what the
original intention could not know. Feedback returns. A sharper intent
appears. The Shot evolves.

Some Shots remain private and serve one human well. Some are published and
invite others inside. Some reach the App Store. Some become communities.
Some acquire appcoins. Some are forgotten. Some survive their creators.

The protocol does not decide which path deserves to happen. Its
responsibility is simpler and harder:

\begin{itemize}
\item keep the factory open;
\item keep the artifact sovereign;
\item keep identity in the human's hands;
\item keep public claims legible;
\item keep the network replaceable;
\item keep history intact;
\item and preserve the connection between an intention and everything
that grows downstream from it.
\end{itemize}

The first stable release will create the \TOH{} mobile app. That app will
become the first canonical Shot of the system that created it. From there,
the circle opens.

And when it does, the invitation will be the same one printed at the end
of every build, the one this whole system exists to make easy to accept:

\vspace{1.5em}
\begin{center}
{\displayfont\Large Take another one.}
\end{center}

\vspace{2em}
\creedrule

% =====================================================================
% COLOPHON
% =====================================================================
\clearpage
\section*{Colophon}

This document is Version 0.2 of the \TOH{} Genesis Whitepaper. It
supersedes the preliminary draft of July 25, 2026, and was written against
factory release 0.5.0, before the first stable release and before the
Genesis Shot exists.

It is versioned deliberately. Version 1.0 of this document is reserved: it
will be finalized alongside the first stable release and signed with the
bootstrap builder identity, so that the constitution and the system it
describes reach stability together. The loom must be complete before it
weaves its first flag.

The document is typeset in EB~Garamond and IBM~Plex~Mono, in the darkroom
palette of the \TOH{} factory: near-black ink, silver-halide grey, and one
hot signal color. Page numbers are exposure numbers. If you are reading
this from a permanent content-addressed store, then this file has become
what it describes: a small signed fact, preserved beyond any single
server, waiting for whoever finds it.

\vspace{1.5em}
{\monolabel\scriptsize\color{halide}
TOHSENO / GENESIS WHITEPAPER / V0.2 / JULY 2026\\[0.3em]
GIVE EVERY IDEA A SHOT. MOST SHOTS MISS. THE PROTOTYPE IS THE PAYOFF.\\[0.3em]
END OF ROLL 01\par}

\end{document}
